\section*{Цель работы}
Экспериментальное исследование условий возникновения резонанса напряжений и токов, изучение амплитудно-частотных характеристик (АЧХ) последовательного и параллельного колебательных контуров, определение влияния параметров цепи на добротность.

\section*{Теоретические сведения}
Резонанс в электрической цепи — режим, при котором входное сопротивление (или проводимость) является чисто активным, а ток и напряжение на входе совпадают по фазе.

\begin{enumerate}
    \item \textit{Последовательный контур (резонанс напряжений).} Условие резонанса: реактивное сопротивление $X = \omega L - \frac{1}{\omega C} = 0$.
    \begin{equation}
        f_0 = \frac{1}{2\pi\sqrt{LC}}
    \end{equation}
    При резонансе полное сопротивление $Z = R$ минимально, а ток $I_0 = \frac{U}{R}$ максимален. Добротность контура $Q$ характеризует превышение напряжения на реактивных элементах над входным:
    \begin{equation}
        Q = \frac{U_{L0}}{U} = \frac{U_{C0}}{U} = \frac{\rho}{R}
    \end{equation}
    где $\rho = \sqrt{\frac{L}{C}}$ — характеристическое сопротивление.

    \item \textit{Параллельный контур (резонанс токов).} Условие резонанса: реактивная проводимость $B = \omega C - \frac{1}{\omega L} = 0$. При резонансе полная проводимость $Y = G$ минимальна, а напряжение $U_0 = \frac{I}{G}$ максимально. Добротность:
    \begin{equation}
        Q = \frac{I_{L0}}{I} = \frac{I_{C0}}{I} = \frac{R}{\rho}
    \end{equation}
\end{enumerate}

АЧХ контура определяется зависимостью модуля проводимости $|Y(j\omega)|$ (для последовательного) или модуля сопротивления $|Z(j\omega)|$ (для параллельного) от частоты. Полоса пропускания $\Delta f$ определяется на уровне $0.707$ от максимума АЧХ, при этом $Q = \frac{f_0}{\Delta f}$.

\section*{Инструкция по выполнению работы}

\subsection*{Исследование резонанса напряжений}
\begin{enumerate}
    \item Соберите схему последовательного контура согласно \cref{fig:series_circuit}.
    \begin{itemize}
        \item Подключите генератор сигналов (ГС) к цепи. Для имитации идеального источника напряжения параллельно выходу ГС подключите низкоомный резистор (согласно указаниям стенда).
        \item Включите в цепь катушку $L$, магазин емкостей $C$ и амперметр.
    \end{itemize}
    \item \textit{Определение резонанса:} Установите напряжение $U = 2$~В. Варьируя частоту ГС в диапазоне $1 \dots 7$~кГц, найдите частоту $f_0$, при которой ток $I$ максимален.
    \item Занесите значения $U, I_0, f_0$ и напряжение на конденсаторе $U_{C0}$ в \cref{tab:series_res}.
    \item \textit{Снятие АЧХ:} Изменяйте частоту от $0.5 f_0$ до $2 f_0$. Снимите 5 точек до резонанса и 5 точек после. Контролируйте постоянство входного напряжения $U = 2$~В. Данные занесите в \cref{tab:series_afc}.
    \item Повторите измерения для контура с введенным дополнительным сопротивлением $R_1$ и при изменении емкости (согласно указаниям преподавателя).
\end{enumerate}

\subsection*{Исследование резонанса токов}
\begin{enumerate}
    \item Соберите схему параллельного контура согласно \cref{fig:parallel_circuit}.
    \begin{itemize}
        \item Для имитации источника тока ГС подключается через высокоомный резистор.
    \end{itemize}
    \item \textit{Определение резонанса:} Установите ток источника $I = 0.5$~мА. Найдите частоту $f_0$, при которой напряжение на контуре $U$ максимально.
    \item Занесите значения $I, U_0, f_0$ и ток в ветви конденсатора $I_{C0}$ в \cref{tab:parallel_res}.
    \item \textit{Снятие АЧХ:} Аналогично последовательному контуру, снимите зависимость $U(f)$ при постоянном токе $I$. Данные занесите в \cref{tab:parallel_afc}.
\end{enumerate}

\section*{Дорожная карта обработки данных}
\begin{enumerate}
    \item По данным резонанса вычислите параметры элементов: $R = \frac{U}{I_0}$, $\rho = \frac{U_{C0}}{I_0}$, $L = \frac{\rho}{2\pi f_0}$, $C = \frac{1}{2\pi f_0 \rho}$.
    \item Рассчитайте экспериментальную добротность двумя способами: $Q_1 = \frac{U_{C0}}{U}$ и $Q_2 = \frac{f_0}{\Delta f}$.
    \item Постройте совмещенные графики АЧХ для разных параметров контура.
\end{enumerate}

\begin{figure}[hbt]
    \centering
    \includegraphics[width=0.6\linewidth]{figs/series_scheme.pdf}
    \caption{Схема исследования резонанса напряжений}
    \label{fig:series_circuit}
\end{figure}

\begin{figure}[hbt]
    \centering
    \includegraphics[width=0.6\linewidth]{figs/parallel_scheme.pdf}
    \caption{Схема исследования резонанса токов}
    \label{fig:parallel_circuit}
\end{figure}

% Соответствие изображений:
% fig:series_circuit -> Рис. 7.3 из методички
% fig:parallel_circuit -> Рис. 7.6 из методички